%% ma: L10-B14-0004
%% lop: 10
%% chuong: 5
%% bai: 14
%% dang: 10.14.2
%% loai: TL
%% muc: VD
%% dap_an: "-"
%% nguon: "chuyen-de-cac-so-dac-trung-cua-mau-so-lieu-khong-ghep-nhom-toan-10-chuong-trinh-moi.pdf (GV. Phan Nhật Linh) (P04), Bài 03 Dạng 1 Bài tập tự luận 4"
%% kien_thuc: "Lớp 10 Bài 13 (tứ phân vị), Bài 14 (khoảng tứ phân vị)"
%% trang_thai: chinh-thuc
%% ghi_chu: "Đề gốc ghi 'học sinh 10A' với 213 học sinh, đã sửa thành học sinh khối 10. Tự duyệt theo yêu cầu giáo viên 22/10/2026."
\begin{ex}
Chiều cao đầu năm học của học sinh khối 10 một trường được cho trong bảng sau:
\begin{center}
\begin{tabular}{|c|c|c|c|c|c|}\hline
Chiều cao (cm) & $150$ & $155$ & $160$ & $165$ & $170$\\\hline
Tần số & $25$ & $28$ & $103$ & $44$ & $13$\\\hline
\end{tabular}
\end{center}
Tìm khoảng biến thiên, khoảng tứ phân vị của mẫu số liệu.
\loigiai{$R=170-150=20$.

$n=213$ (lẻ): $Q_2=x_{107}=160$ (tần số tích lũy $25$; $53$; $156$; $200$; $213$).

Nửa dưới gồm $106$ giá trị đầu: $Q_1=\dfrac{x_{53}+x_{54}}2=\dfrac{155+160}2=157{,}5$.

Nửa trên gồm $106$ giá trị cuối ($x_{108}$ đến $x_{213}$): $Q_3=\dfrac{x_{160}+x_{161}}2=165$.

$\Delta_Q=165-157{,}5=7{,}5$.}
\end{ex}
